\chapter{Evaluation Setup}
This chapter defines the framework for assessing the effectiveness and efficiency of the proposed semantic enrichment and retrieval system. We outline the evaluation objectives, the process of constructing a "gold standard" dataset, the metrics used for quantitative analysis, and the suite of baselines and configurations used for comparative study.

\section{Evaluation Objectives}
The evaluation is designed to validate the system's ability to structure noisy internship data accurately while adhering to the performance constraints of a lightweight, low-cost application. We formulate our assessment around four primary Research Questions (RQs):

\begin{itemize}
    \item \textbf{RQ1 (Extraction Fidelity):} How accurately does the deterministic pipeline (Layer A) extract structured fields such as compensation, locations, and dates compared to human-annotated labels?
    \item \textbf{RQ2 (Semantic Alignment):} To what extent does the embedding-based mapping (Layer B) improve the classification of job titles into canonical ESCO occupations over keyword-based heuristics?
    \item \textbf{RQ3 (Retrieval Quality):} Does the hybrid search strategy—combining lexical BM25 scores with semantic vector similarity—significantly improve precision and relevance for student-centric queries?
    \item \textbf{RQ4 (System Efficiency):} What are the computational costs (latency, memory, and API overhead) of the multi-layered pipeline, and how does it scale on commodity hardware compared to LLM-only approaches?
\end{itemize}

\section{Ground Truth Construction}
To provide a reliable benchmark, we curated a "Gold Standard" evaluation set. This dataset serves as the reference point for all accuracy and retrieval metrics.

\subsection{Dataset Composition}
The evaluation set consists of 500 internship listings, purposively sampled to ensure diversity across sources and domains:
\begin{itemize}
    \item \textbf{Web-Scraped (400):} Randomly sampled from the 17,500+ records in the primary JSON catalog, covering various platforms like LinkedIn and Indeed.
    \item \textbf{User-Uploaded (100):} Selected from raw JSON uploads to test the system's robustness against heterogeneous, non-standardized structures.
\end{itemize}

\subsection{Annotation Process}
Two domain experts (graduate students in Computer Science) manually labeled each listing according to the target schema defined in Chapter 7. For each entry, annotators:
\begin{enumerate}
    \item Extracted ground-truth city, country, and compensation ranges.
    \item Mapped the raw job title to the most specific "Level 4" ESCO Occupation.
    \item Identified core technical skills from the ESCO Skills pillar.
\end{enumerate}
Inter-annotator agreement was measured using Cohen's Kappa ($\kappa$), achieving a score of 0.82, which indicates high reliability. Discrepancies were resolved through a final review by a third senior annotator.

\section{Metrics}
We utilize a multi-faceted set of metrics to capture different aspects of system performance.

\subsection{Extraction and Classification}
For field-level extraction (Layer A) and categorical mapping (Layer B), we use standard classification metrics:
\begin{itemize}
    \item \textbf{Precision (P):} The proportion of extracted entities that are correct.
    \item \textbf{Recall (R):} The proportion of ground truth entities correctly extracted.
    \item \textbf{F1-Score:} The harmonic mean of precision and recall: $F1 = 2 \cdot \frac{P \cdot R}{P + R}$.
    \item \textbf{Top-k Accuracy:} For title mapping, we report Top-1 and Top-3 accuracy to account for valid overlapping professional families in the ESCO taxonomy.
\end{itemize}

\subsection{Information Retrieval (IR)}
For search engine evaluation, we use 20 representative search queries. Annotators marked each result in the Top-10 as "Highly Relevant," "Relevant," or "Irrelevant."
\begin{itemize}
    \item \textbf{Precision@k (P@k):} Precision at rank $k$ (typically $k=5, 10$).
    \item \textbf{Mean Average Precision (MAP):} A global measure of ranking quality across all queries.
    \item \textbf{NDCG (Normalized Discounted Cumulative Gain):} Accounts for graded relevance levels, rewarding systems that place "Highly Relevant" results higher.
\end{itemize}

\section{Experimental Configuration}
To isolate the contribution of each component, we compare our proposed system against several baselines and ablation variants.

\subsection{Baselines}
\begin{itemize}
    \item \textbf{Baseline 1: Keyword-Only (BM25):} A traditional search on raw text without any semantic enrichment.
    \item \textbf{Baseline 2: Rules-Only (Layer A):} Extraction using only regex and lexicons, representing a basic structured system.
    \item \textbf{Baseline 3: Hybrid (Layer A + Layer B):} Our primary proposed system, combining deterministic rules with small-model semantic mapping.
\end{itemize}

\subsection{Comparison and Hardware}
To evaluate efficiency (RQ4), we compare the throughput and cost of the local \textbf{Layer B} (Transformer-based) against an \textbf{LLM-Only (Layer C)} approach using GPT-4o-mini.

All local tests are conducted on a standardized environment:
\begin{itemize}
    \item \textbf{CPU:} Intel Core i7 (4-core), 16GB RAM (No dedicated GPU).
    \item \textbf{Runtime:} Node.js v20 (Next.js 14 Production Build).
    \item \textbf{Embedding Model:} \texttt{all-MiniLM-L6-v2} (ONNX format) for local inference.
\end{itemize}
